\labsection{6}{Build and extend a student data frame}{sec:lab06}

\begin{labproject}{Lab 6 guided practice}
\textbf{Coverage.} Chapter~\ref{chap:data-frames}.\\
\textbf{Goal.} Build and extend a student table, then inspect its structure and summaries.

\textbf{Task 1 --- Create the base data frame.}

\begin{lstlisting}[style=dsr]
students <- data.frame(
  name = c(
    "Anastasia", "Dima", "Michael", "Matthew",
    "Laura", "Kevin", "Jonas"
  ),
  score = c(12.5, 9.0, 16.5, 12.0, 9.0, 8.0, 19.0),
  attempts = c(1, 3, 2, 3, 2, 1, 2)
)

print(students)
\end{lstlisting}

\textbf{Task 2 --- Add the specified \texttt{qualify} column.}

\begin{lstlisting}[style=dsr]
students$qualify <- c(
  "yes", "no", "yes", "no", "no", "no", "yes"
)
print(students)
\end{lstlisting}

\textbf{Task 3 --- Add Emily's row.}

\begin{lstlisting}[style=dsr]
emily <- data.frame(
  name = "Emily",
  score = 14.5,
  attempts = 1,
  qualify = "yes"
)

students <- rbind(students, emily)
print(students)
\end{lstlisting}

\textbf{Task 4 --- Inspect the dataset.}

\begin{lstlisting}[style=dsr]
str(students)
summary(students)
dim(students)
nrow(students)
ncol(students)
\end{lstlisting}

Convert \texttt{qualify} to a factor to see how its summary changes:

\begin{lstlisting}[style=dsr]
students$qualify <- factor(students$qualify)
summary(students)
\end{lstlisting}

The final table has 8 rows, 4 columns, and scores from 8.0 to 19.0.
\end{labproject}

\begin{labdeliverable}
Submit the table before and after modification, its structure and summaries, dimensions, and a brief note on the factor conversion.
\end{labdeliverable}
